\chapter{Getting started}
\label{ch:start}

\section{Building}

The build is that of GROMACS 4.5.4 with CMake; FFTW~3 in single precision is needed for the electron cloud.
From the source tree:
\begin{lstlisting}
cd gromacs-4.5.4-SFMD
mkdir build && cd build
cmake .. -DCMAKE_BUILD_TYPE=Release -DGMX_FFT_LIBRARY=fftw3 \
         -DCMAKE_INSTALL_PREFIX=/path/to/install
make -j8 install
\end{lstlisting}
For a reference build that must not change when libraries are rebuilt, add \code{-DBUILD\_SHARED\_LIBS=OFF}.
Run inputs must be written with this build's \code{grompp}: the \file{.tpr} carries the \opt{sfmd-*} options (format
185), which stock GROMACS and \code{gromacs-mc} cannot read.

\section{A first run}

\code{example/} is a complete run: lysozyme (PDB 1AKI, 1960 atoms) in vacuum, exploded by a 5\,mJ, 5\,fs pulse.
The steps are those of any GROMACS run, with \code{-ionize} added to \code{mdrun}:
\begin{lstlisting}
pdb2gmx -f 1aki.pdb -ff charmm27 -water none
grompp  -f exp.mdp -c conf.gro -p topol.top -o explode.tpr
mdrun   -deffnm explode -nt 4 -pd -ionize
\end{lstlisting}
\code{-water none} keeps \code{pdb2gmx} from asking; \code{-pd} (particle decomposition) is required for more than one
thread, because domain decomposition cannot follow an explosion (Appendix~\ref{app:parallel}).

The \file{.mdp} needs, besides the usual vacuum settings,
\begin{lstlisting}
dt                          = 0.000002   ; 2 as (ps!)
morse                       = yes        ; bonds that can break
pbc                         = no
coulombtype                 = Cut-off
sfmd-pulse-peak-time        = 35         ; fs
sfmd-pulse-fwhm             = 15         ; fs, full width at half maximum
sfmd-pulse-energy           = 1.13e-4    ; J
sfmd-pulse-focal-diameter   = 30000      ; nm
sfmd-pulse-wavelength       = 850        ; nm
\end{lstlisting}
Note the units: pulse times in fs, \code{dt} and \code{tinit} in GROMACS's ps.
These five pulse values give a peak intensity of $10^{15}$\,W/cm$^2$; the run summary prints it
(Section~\ref{sec:summaries}).
The cut-offs are set by \code{grompp} when the cloud is on (\opt{sfmd-pb-long-range = auto},
Section~\ref{sec:longrange}); without the cloud use \code{rcoulomb = rvdw = rlist = 0} (all pairs) or a cut-off that
covers the exploding sample.

\section{With the trapped-electron cloud}

\begin{lstlisting}
sfmd-electrons          = pbgrid
sfmd-coll-recombination = 2      ; optional: recombination after the pulse
\end{lstlisting}
With \opt{sfmd-electrons = pbgrid} the cloud's options default to the recommended model: the laser barrier, laser
heating of an ellipsoidal cloud with surface damping, electron-impact ionization and a 1\,eV temperature floor
(Chapters~\ref{ch:pb}--\ref{ch:coll}).
Each can be switched off on its own.
Recombination stays off by default; after the pulse (\opt{= 2}) it gives back most of the extra charge impact
ionization creates without changing the explosion.

\devnote{Before 11 November 2026 the cloud's options all defaulted to 0. An \file{.mdp} written then that relied on the
defaults now gets the full model; set the options to 0 (and \opt{sfmd-heating-shape = sphere}) to reproduce it. A
\file{.tpr} keeps what it was made with.}

\section{How long to run}

The ions' directions settle within roughly 100\,fs of the pulse; their kinetic energy keeps growing slowly much longer.
Two options cut the cost:
\begin{description}
  \item[\opt{sfmd-autostop = 2}] stops the run once the heavy-ion velocity directions have settled: every
    \opt{sfmd-autostop-interval} (5\,fs) the median turn since the last check is extrapolated as $1/(t-t_{\mathrm{peak}})^2$,
    and the run stops when the turn still to come is below \opt{sfmd-autostop-angle} (2$^\circ$), but not before
    \opt{sfmd-autostop-min-time} (60\,fs) after the pulse peak; \code{nsteps} stays the upper limit.
    On lysozyme it stopped at 95\,fs ($10^{16}$) and 125--130\,fs ($10^{15}$), with the directions still turning
    0.4--2.2$^\circ$ by 300\,fs against a seed spread of 1--5$^\circ$, and the ion energy 6\,\% short of its 300\,fs value;
    2.5--3.3$\times$ less wall time.
    \opt{sfmd-autostop = 1} is the older stop at a ratio of kinetic to total energy.
  \item[\opt{sfmd-schedule}] lengthens the step and the cloud-solve interval after the pulse, e.g.
    \code{65:0.01:0.2 95:0.1:2 135:0.1:5} (from 65\,fs a 10\,as step and a solve every 0.2\,fs, and so on).
    \code{nsteps} $\times$ \code{dt} stays the end time.
    \code{grompp} refuses a first change before the pulse peak + $3\sigma$ (changing the step while ADK still ionizes
    changes the run's random history) and output intervals the new steps do not divide.
    On lysozyme the schedule above moved the directions by at most 1.2$^\circ$ and the ion energy by $-0.2$\,\% at
    300\,fs, 2.8--3.1$\times$ faster.
    Off by default.
\end{description}
Compare kinetic energies only between runs stopped or scheduled the same way.

\section{Seeds and ensembles}
\label{sec:seeds}

The ADK draws are seeded from \code{/dev/random}, so every run is an independent realization --- what an ensemble
needs.
\env{GMX\_SFMD\_SEED=n} pins the seed; each rank still gets its own stream (the rank index is mixed in), so a pinned
run is reproducible for a given rank count only.
The seed used is in the run summary.
Fixed seeds are for tests, not for production ensembles.

\section{Chaining runs}

\opt{sfmd-initial-charges = 1} starts from the charges in \file{charges.txt} (atom index and charge per line), which a
finished run writes into \file{simulation\_output/}.
It cannot be used with the cloud, whose state is not saved; \code{grompp} refuses the combination.

\section{The run summary}

At the first step \code{mdrun} prints a summary of what the run will do --- system, pulse with its peak intensity and
field, length, schedule, seed, ranks, every process with its settings --- and warnings for settings outside what was
tested (Section~\ref{sec:summaries}).
Read it first.
